% Данный файл распространяетсяя под лицензией CC BY 4.0. Текст лицензии размещён на https://creativecommons.org/licenses/by/4.0/
% (c) Кафедра системного программирования, 2025

\documentclass[a4paper]{article}

\usepackage[a4paper, top=8mm, bottom=15mm, left=8mm, right=8mm]{geometry}

\usepackage{polyglossia}
\setdefaultlanguage[babelshorthands=true]{russian}
\setotherlanguage{english}

\usepackage{fontspec}
\setmainfont{FreeSerif}
\newfontfamily{\russianfonttt}[Scale=0.7]{DejaVuSansMono}

\usepackage[tiny, compact]{titlesec}

\usepackage{titling}
\setlength{\droptitle}{-1cm}
\pretitle{\begin{center}\begin{bfseries}\Large}
\posttitle{\par\end{bfseries}\end{center}}
\preauthor{\begin{center}\normalsize}
\postauthor{\par\end{center}\vspace{-1.8cm}}

\usepackage{hyperref}
\usepackage{bookmark}
\usepackage{csquotes}
\usepackage{fancyhdr}
\setlength{\footskip}{8mm}
\fancyhf{}
\fancyfoot[C]{\footnotesize Moroz Vlasdislav SPbu 2026}
\renewcommand{\headrulewidth}{0pt}
\renewcommand{\footrulewidth}{0pt}
\pagestyle{fancy}
\fancypagestyle{plain}{%
    \fancyhf{}
    \fancyfoot[C]{\footnotesize Moroz Vlasdislav SPbu 2026}
    \renewcommand{\headrulewidth}{0pt}
    \renewcommand{\footrulewidth}{0pt}
}
% Сюда можно дописывать нужные вам пакеты.

\title{Разработка обратного индекса для базовых чанков в алгоритмах SBC}

\author{Мороз Владислав Владимирович}

% Дата много места занимает, её лучше не указывать.
\date{}

\begin{document}

\maketitle

\begin{flushright}
    Группа: \emph{группа в формате \enquote{25.Б72-мм}}

    Кафедра: \emph{системного программирования}

    % Степень/звание и должность научника мы лучше вас знаем, их можно не указывать.
    Научный руководитель: \emph{Гориховский Вячеслав Игоревич}

    % А вот про консультанта укажите официальное место работы, с "ООО" или чем-то таким, и должностью желательно по трудовой, не просто "techlead". Выясните у консультанта.
    Консультант: \emph{Дмитриевцев Алексей Сергеевич, младший инженер по разработке ПО, ООО "КНС ГРУПП"}
    
    Номер семестра практики: \emph{3}
\end{flushright} 


\section{Постановка задачи}

Системы дедупликации уменьшают объём хранимых данных, выявляя совпадающие или похожие фрагменты. Для разбиения данных на чанки применяются, в частности, алгоритмы Content-Defined Chunking (CDC), а для поиска похожих чанков и повторного использования их содержимого --- алгоритмы Similarity-Based Chunking (SBC) \cite{xia2020,aronovich2009}. В SBC каждый чанк описывается компактным представлением (sketch), по которому оценивается сходство с другими чанками. Подходы к быстрому обнаружению похожих объектов, включая Finesse и ODESS, также используют признаки чанков для сокращения поиска по хранилищу \cite{zhang2019,zou2021}.

Задача поиска базового чанка состоит в том, чтобы для нового sketch найти уже сохранённые чанки с достаточно высокой степенью сходства. При прямом переборе для каждого запроса приходится сравнивать его со всеми сохранёнными чанками; по мере роста хранилища число таких сравнений становится существенным. Поэтому требуется структура, которая быстро отбирает небольшое множество кандидатов, не исключая преждевременно потенциально подходящие чанки.

\textbf{Постановка задачи.} Требуется разработать обратный индекс базовых чанков для алгоритмов SBC. Индекс должен хранить sketch чанков и обратные списки (\emph{posting lists}), связывающие признак с идентификаторами чанков, в которых этот признак встречается. По запросу индекс должен формировать множество кандидатов по общим признакам, после чего позволять ранжировать их по точной степени сходства. При этом необходимо поддержать согласованное изменение записей индекса и исследовать компромисс между сокращением числа сравнений и качеством поиска.

\textbf{Цель работы} --- разработать и исследовать эффективный обратный индекс для поиска базовых чанков в алгоритмах SBC, сокращающий число сравнений при поиске похожих объектов и пригодный для интеграции в pipeline дедупликации библиотеки MarLine. При оценке решения необходимо учитывать производительность индекса, потребление памяти и влияние параметров поиска на скорость и качество результатов.

Для достижения цели необходимо решить следующие задачи:
\begin{enumerate}
    \item Спроектировать архитектуру и программный интерфейс индекса для хранения идентификаторов чанков, их sketch-представлений и признаков.
    \item Разработать структуру хранения sketch и posting lists, включая операции чтения, добавления, удаления и очистки записей.
    \item Обеспечить согласованное добавление, удаление и обновление записей, в том числе определить поведение при повторных вставках.
    \item Разработать поиск top-$k$ наиболее похожих чанков: формировать кандидатов по общим признакам и точно оценивать только отобранные объекты.
    \item Реализовать вычисление коэффициента Жаккара для sketch различного размера и с различными типами признаков.
    \item Исследовать эвристики сокращения множества кандидатов: ограничение частоты признаков, приоритет редких признаков и отсечение по порогу сходства.
    \item Реализовать сбор статистики ложноположительных результатов и исследовать возможность использовать её для адаптации поиска.
    \item Оценить влияние индекса и параметров поиска на число рассматриваемых кандидатов, время работы, потребление памяти и качество результатов.
    \item Подготовить модульные тесты, проверяющие соответствие индекса и операций с ним заданной спецификации.
\end{enumerate}

\section{Описание предлагаемого решения}

% Раздел пишется в свободной форме, с минимумом технических подробностей (можно приводить ссылки на более подробные документы).

Тут надо кратко описать идею решения, спозиционировать его относительно предыдущих наработок\footnote{Если на них надо сослаться, то в подстраничной сноске, URL: \url{http://example.com} (дата обращения: не забываем указывать).} и существующих аналогов (тоже предельно кратко), по возможности обосновать принятые решения.

% Вот это важно, опишите кратко, что использовали
Реализация решения осуществлялась с использованием \dots (тут описать используемые технологии).

\section{Апробация / Эксперименты}

% Название секции можно поменять на более подходящее вашей работе

Название этой секции --- либо \enquote{Апробация}, либо \enquote{эксперименты} (если содержательно есть и то и то, можно две отдельные секции сделать). Тут опишите, как проводилась проверка того, что получилось что-то адекватное. Если есть отзывы пользователей, отлично (особенно если они в структурированном виде, типа анкеты SUS). Если есть отзывы команды, консультанта и т.п., хорошо. Если есть только тесты. то хотя бы их тут опишите.

Если работа подразумевала замеры чего-то (например, производительности), приведите тут итоговые результаты и выводы, и дайте ссылку на таблицы с данными. Не забывайте про матстат (приводите матожидание и среднеквадратичное отклонение, не просто результат замера).

Приведите выводы из эксперимента.

\section{Заключение}

В ходе выполнения работы получены следующие результаты:

% ниже списком перечисляете то, что выносится как результат практики.

\begin{itemize}
    \item \dots
    \item \dots
    \item \dots
\end{itemize}

Материалы, иллюстрирующие выполнение работы, размещены по адресу: \url{http://www.example.com} (если есть и применимо --- скриншоты, видео на 1-2 минуты и т.п. с демонстрацией; может быть очень полезно даже для консольных утилит или библиотек).

Репозиторий/ссылка на пуллреквест: \url{http://www.example.com} (это обязательно; ссылок может быть несколько).

Техническая документация: \url{http://www.example.com} (если применимо --- для пуллреквестов можно техническую сторону прямо в комментарии к пуллреквесту указать, но можно и в отдельную вики-страницу вынести).

Решение внедрено в производственные процессы компании/прошло процесс ревью/находится на ревью/находится в разработке и т.п. В общем, кратко сформулируйте, что в итоге получилось и в каком статусе результат, чтобы не было ощущения, что оно сделано, положено на GitHub и забыто.

\begin{thebibliography}{9}
    \bibitem{xia2020}
    Xia W., Zou X., Jiang H., Zhou Y., Zhang C. et al.
    \newblock The Design of Fast Content-Defined Chunking for Data Deduplication-Based Storage Systems.
    \newblock IEEE Transactions on Parallel and Distributed Systems. 2020. Vol.~31, no.~9. P.~2017--2031.
    \newblock DOI: \href{https://doi.org/10.1109/TPDS.2020.2984632}{10.1109/TPDS.2020.2984632}.

    \bibitem{aronovich2009}
    Aronovich L., Asher R., Bachmat E., Bitner H., Hirsch M., Klein S.~T.
    \newblock The Design of a Similarity Based Deduplication System.
    \newblock In: Proceedings of the Israeli Experimental Systems Conference (SYSTOR 2009). 2009.
    \newblock DOI: \href{https://doi.org/10.1145/1534530.1534539}{10.1145/1534530.1534539}.

    \bibitem{zhang2019}
    Zhang Y., Xia W., Feng D., Jiang H., Hua Y., Wang Q.
    \newblock Finesse: Fine-Grained Feature Locality Based Fast Resemblance Detection for Post-Deduplication Delta Compression.
    \newblock In: Proceedings of the 17th USENIX Conference on File and Storage Technologies (FAST 2019). 2019. P.~121--128.
    \newblock URL: \url{https://www.usenix.org/conference/fast19/presentation/zhang}.

    \bibitem{zou2021}
    Zou X., Deng C., Xia W., Shilane P., Tan H., Zhang H., Wang X.
    \newblock ODESS: Speeding up Resemblance Detection for Redundancy Elimination by Fast Content-Defined Sampling.
    \newblock In: Proceedings of the 37th IEEE International Conference on Data Engineering (ICDE 2021). 2021.
    \newblock DOI: \href{https://doi.org/10.1109/ICDE51399.2021.00048}{10.1109/ICDE51399.2021.00048}.
\end{thebibliography}

\end{document}
